\section{Introduction}
\label{sec:introduction}

Article~1 of this thesis showed that automatically inferred JML specifications materially improve LLM-generated unit tests on Apache Commons Lang. The empirical claim assumed that the consumer of the specifications --- the test generator, an IDE, a verifier --- has access to the same source tree the inferrer ran over. That assumption is rarely satisfied in practice. Java libraries are distributed as JAR files; REST endpoints expose only their HTTP signature; gRPC services expose only their protobuf descriptors. Source code is the exception, not the norm.

This article asks: \emph{can inferred specifications be embedded in compiled artefacts and interface descriptions, with sufficient fidelity that downstream consumers can use them as if they had the source?} The empirical question is operational. The conceptual question --- whether specifications \emph{should} be transported alongside compiled artefacts --- has a long history~\cite{meyer1992design,leavens2006design,jml}; the practical question is whether the standard toolchain can be made to carry them without modification. The answer in this article is yes, with 100.0\% fidelity on three real-world libraries (Section~\ref{sec:empirical}) and an 11.76\%--13.49\% byte overhead.

The contribution is the embedder/extractor pair plus its OpenAPI counterpart. Three concrete artefacts:

\begin{enumerate}
    \item A Java annotation type \texttt{@JmlSpec}, repeatable per clause, written by an ASM-based embedder into class-file \texttt{RuntimeVisibleAnnotations} attributes. Reflection or any class-file reader recovers the embedded specification.
    \item An OpenAPI~3.x extension family (\texttt{x-jml-requires}, \texttt{x-jml-ensures}, \texttt{x-jml-assignable}, \texttt{x-jml-signals}, \texttt{x-jml-invariant}) carried per-operation and per-schema. The extension is preserved by standard OpenAPI parsers (swagger-parser, openapi-generator) without modification.
    \item A Maven classifier sidecar (\texttt{<artifact>-jml.jar}) carrying JML stub files, used when in-bytecode embedding is unsuitable (signed JARs that cannot be re-signed; specifications too large to fit the per-class budget).
\end{enumerate}

The article evaluates the bytecode side empirically and the OpenAPI side by construction. The principal empirical findings:

\begin{itemize}
    \item Roundtrip fidelity is 100.0\% on all three libraries tested (20{,}546 methods total) plus the real-inferred-specs corpus, well above the 95\% threshold required for the system to be useful in CI/CD.
    \item Bytecode-size overhead is in the 11.76\%--13.49\% range under a deliberately bulky synthetic-spec workload. Real inferred specifications are typically smaller per method, so this is an upper bound.
    \item Embedding and reading are fast enough --- in the 10{,}000--235{,}000 methods/second range --- to support per-pull-request CI workflows on libraries an order of magnitude larger than those tested.
    \item Classes embedded with \texttt{@JmlSpec} load correctly in JVMs with no awareness of the annotation type. The documented constraint is that reflective annotation access may raise \texttt{TypeNotPresentException}; class load itself succeeds.
\end{itemize}

The remainder of the article is organised as follows. Section~\ref{sec:background} reviews the existing channels for shipping specifications alongside compiled Java code, and the gap each of them leaves. Section~\ref{sec:format} presents the format design --- the annotation type, the canonical-form printer, and the OpenAPI extension schema. Section~\ref{sec:impl} describes the implementation. Section~\ref{sec:empirical} reports the empirical validation. Section~\ref{sec:discussion} discusses generalisability; Section~\ref{sec:threats} the threats to validity. Section~\ref{sec:related} surveys related work and Section~\ref{sec:conclusion} concludes.
